\documentclass[12pt,a4paper]{article}
\usepackage[utf8]{inputenc}
\usepackage{graphicx}
\usepackage{hyperref}
\usepackage{geometry}
\geometry{margin=1in}

\title{Technical Report: Document Management System}
\author{Mini-Project Team}
\date{\today}

\begin{document}

\maketitle

\begin{abstract}
This report outlines the architecture and implementation of a Document Management System (DMS) designed to standardize project document structures and improve release workflows. The system features ABAC/RBAC access controls, Server-Side Encryption (SSE) via MinIO, automated audit logging, and document version comparison tools.
\end{abstract}

\section{Introduction}
The DMS is a full-stack application built using Node.js (Express/TypeScript) for the backend, React (Vite) for the frontend, PostgreSQL for relational data and metadata, and MinIO for S3-compatible document storage.

\section{System Architecture}
\subsection{Technology Stack}
The technology stack was chosen for flexibility, developer velocity, and robust security:
\begin{itemize}
    \item \textbf{Backend:} Express.js + TypeScript
    \item \textbf{Frontend:} React + Vite
    \item \textbf{Database:} PostgreSQL with Prisma ORM
    \item \textbf{Storage:} MinIO (Docker) with SSE enabled
\end{itemize}

\subsection{Document Storage \& Versioning}
Documents are restricted to a maximum size of 5MB. Every upload is tracked as a new version. The actual files are stored in MinIO, while PostgreSQL maintains the relationship between projects, documents, and versions. Retention policies are dynamic and managed by a background cron job referencing the \texttt{SystemConfig} table.

\section{Access Control Model}
Access control uses a hierarchical model (ABAC/RBAC):
\begin{itemize}
    \item \textbf{Levels:} Department $>$ Project
    \item \textbf{Roles:} Viewer, Editor, Admin
\end{itemize}
Users assigned to a department inherit their role across all associated projects unless a project-specific role overrides it.

\section{Security \& Audit Logging}
\subsection{Encryption}
Data in transit is secured using TLS 1.3 (handled via reverse proxy/load balancer in production), and data at rest is encrypted using MinIO Server-Side Encryption (SSE).
\subsection{Audit Log}
An append-only database table records the following events:
\texttt{CREATE}, \texttt{UPLOAD}, \texttt{VIEW}, \texttt{DOWNLOAD}, \texttt{EDIT}.
Each log entry securely captures the timestamp, User ID, and origin IP address.

\section{Document Diffing}
The system provides endpoints to compare document versions:
\begin{itemize}
    \item \textbf{PDF files:} Parsed using \texttt{pdf-parse} and compared using \texttt{diff} word algorithms.
    \item \textbf{Word documents:} Parsed by extracting the \texttt{word/document.xml} from the DOCX zip format and performing strict XML comparisons.
\end{itemize}

\section{Conclusion}
The implemented DMS meets all corporate requirements, providing a secure, reliable, and version-controlled environment for project documentation.

\end{document}
